% ============================================================
%  RAPPORT D'AVANCEMENT DE THÈSE — DOCTORAT INFORMATIQUE
%  Federated Time Series Foundation Models
%  pour la Maintenance Prédictive Industrielle Multi-Domaine
% ============================================================
\documentclass[12pt, a4paper]{report}

% --- Packages ---
\usepackage[utf8]{inputenc}
\usepackage[T1]{fontenc}
\usepackage[french]{babel}
\usepackage[margin=2.5cm]{geometry}
\usepackage{graphicx}
\usepackage{booktabs}
\usepackage{array}
\usepackage{longtable}
\usepackage{xcolor}
\usepackage{colortbl}
\usepackage{hyperref}
\usepackage{enumitem}
\usepackage{fancyhdr}
\usepackage{titlesec}
\usepackage{parskip}
\usepackage{mdframed}
\usepackage{multirow}
\usepackage{tabularx}
\usepackage{tikz}
\usetikzlibrary{positioning}
\usepackage{caption}

% --- Colors ---
\definecolor{mainblue}{RGB}{30, 80, 160}
\definecolor{lightblue}{RGB}{220, 235, 255}
\definecolor{accentorange}{RGB}{220, 100, 30}
\definecolor{gapred}{RGB}{200, 50, 50}
\definecolor{gaporange}{RGB}{220, 140, 30}
\definecolor{gapyellow}{RGB}{180, 160, 30}
\definecolor{lightgray}{RGB}{245, 245, 245}
\definecolor{darkgray}{RGB}{80, 80, 80}

% --- Hyperref setup ---
\hypersetup{
    colorlinks=true,
    linkcolor=mainblue,
    urlcolor=mainblue,
    citecolor=darkgray,
    pdftitle={Rapport d'avancement de thèse - Yassire AMMOURI},
    pdfauthor={Yassire AMMOURI},
}

% --- Section formatting ---
\titleformat{\chapter}[hang]
  {\huge\bfseries\color{mainblue}}
  {\thechapter.}{12pt}{}
\titleformat{\section}[hang]
  {\Large\bfseries\color{mainblue}}
  {\thesection}{10pt}{}
\titleformat{\subsection}[hang]
  {\large\bfseries\color{darkgray}}
  {\thesubsection}{8pt}{}

% --- Header / Footer ---
\pagestyle{fancy}
\fancyhf{}
\fancyhead[L]{\small\color{darkgray}Université Mohammed V — LRIT}
\fancyhead[R]{\small\color{darkgray}\leftmark}
\fancyfoot[C]{\small\color{darkgray}\thepage}
\fancyfoot[R]{\small\color{darkgray}Rapport d'avancement de thèse}
\fancyfoot[L]{\small\color{darkgray}Y. AMMOURI}
\renewcommand{\headrulewidth}{0.4pt}

% --- Box environment (transparent style) ---
\newmdenv[
  backgroundcolor=white,
  linecolor=darkgray,
  linewidth=0.5pt,
  roundcorner=3pt,
  innertopmargin=6pt,
  innerbottommargin=6pt,
  innerleftmargin=10pt,
  innerrightmargin=10pt,
]{infobox}

\newmdenv[
  backgroundcolor=white,
  linecolor=darkgray,
  linewidth=0.5pt,
  roundcorner=3pt,
  innertopmargin=6pt,
  innerbottommargin=6pt,
  innerleftmargin=10pt,
  innerrightmargin=10pt,
]{notebox}

% ============================================================
\begin{document}

% --- Title Page ---
\begin{titlepage}
  \centering
  \vspace*{1.5cm}
  {\color{mainblue}\rule{\linewidth}{2pt}}\\[0.4cm]
  {\huge\bfseries\color{mainblue} Rapport d'avancement de thèse}\\[0.3cm]
  {\large\color{darkgray} Première année de doctorat}\\[0.2cm]
  {\color{mainblue}\rule{\linewidth}{2pt}}\\[1.2cm]

  {\LARGE\bfseries Federated Time Series Foundation Models\\[0.3cm]
  pour la Maintenance Prédictive Industrielle\\[0.3cm]
  Multi-Domaine}\\[2cm]

  \begin{tabular}{rl}
    \textbf{Préparé par :} & Yassire AMMOURI\\[4pt]
    \textbf{Directeur de thèse :} & Pr. Younes Karfa Bakali\\[4pt]
    \textbf{Encadrant :} & Mme Rajaa SAIDI\\[4pt]
    \textbf{Université :} & Université Mohammed V, Faculté des Sciences de Rabat\\[4pt]
    \textbf{Laboratoire :} & LRIT (Laboratoire de Recherche en Informatique et Télécommunications)\\[4pt]
    % \textbf{Financement :} & Bourse Doctorant Moniteur PASS\\[4pt]
    % \textbf{Date de début :} & Octobre 2025\\[4pt]
    \textbf{Date :}        & \today \\[4pt]
    % \textbf{Destinataires :} & Directeurs de thèse\\
  \end{tabular}

  \vfill
  {\color{darkgray}\small Ce rapport couvre l'ensemble des activités réalisées depuis le
  début de la thèse, incluant le bilan personnel, la revue bibliographique,
  l'identification des problématiques, et la planification des contributions.}
\end{titlepage}

% --- Table of Contents ---
\tableofcontents
\clearpage

% ============================================================
% PART I — ACTIVITÉS RÉCENTES ET PROGRÈS PERSONNELS
% ============================================================
\chapter*{Avant-propos : Activités récentes et progrès}
\addcontentsline{toc}{chapter}{Avant-propos : Activités récentes et progrès}
\markboth{Avant-propos}{}

Ce premier chapitre présente un bilan honnête de ce qui a été accompli depuis le début de
la thèse. Il ne s'agit pas d'un état des lieux de la recherche, mais d'un inventaire personnel :
les formations suivies, les articles en cours, et les responsabilités pédagogiques assumées.
Ce contexte est indispensable pour que les directeurs de thèse puissent évaluer la charge
de travail réelle et guider les priorités des prochains mois.

% -------------------------------------------------------
\section*{1. Formations et certifications suivies}
\addcontentsline{toc}{section}{1. Formations et certifications suivies}

\begin{table}[h!]
\centering
\caption{Certifications obtenues en première année}
\label{tab:certifications}
\renewcommand{\arraystretch}{1.4}
\begin{tabularx}{\textwidth}{>{\bfseries}lXr}
\toprule
Organisme & Description de la formation & Progressent\\
\midrule
WIPO — DL101 & Introduction à la propriété intellectuelle. Formation en ligne de l'Organisation
Mondiale de la Propriété Intellectuelle couvrant les brevets, droits d'auteur, marques
et secrets commerciaux. Compétence utile pour protéger les futures contributions
algorithmiques ou logicielles de la thèse. & complète \\[4pt]
Wiley & Formation à la publication scientifique : comprendre le processus de soumission,
la révision par les pairs, et les bonnes pratiques d'écriture académique. & 1 Certif \\[4pt]
Scopus (Elsevier) & Maîtrise de la base de données bibliographique Scopus : recherche
avancée, analyse de citations, identification des journaux et des tendances dans un
domaine. &  1 Certif \\[4pt]
Springer Nature & Formation à l'utilisation de la plateforme Springer pour la recherche
bibliographique et la compréhension des politiques de publication en accès ouvert. & 1 Certif \\[4pt]
Taylor \& Francis & Formation aux outils de recherche et de publication du groupe
Taylor \& Francis, incluant les journaux spécialisés en ingénierie et en informatique. & 5 Certif \\[4pt]
Clarivate (Web of Science) & Utilisation de Web of Science pour les recherches bibliographiques,
l'analyse des facteurs d'impact et le suivi des indicateurs bibliométriques. & 10 Certif \\
\bottomrule
\end{tabularx}
\end{table}

Au cours de cette première période de thèse, un effort important a été fourni pour
se former aux outils et aux pratiques de la recherche scientifique. Les formations
ci-dessous ont toutes été complétées :

\begin{infobox}
\textbf{Ce que ces formations apportent concrètement :} maîtriser ces outils signifie
être capable de conduire une revue de littérature rigoureuse, identifier les bons journaux
pour soumettre, et comprendre les mécanismes de la propriété intellectuelle qui entourent
les contributions de recherche. C'est un socle essentiel pour toute la suite du doctorat.
\end{infobox}

% -------------------------------------------------------
\section*{2. Avancement de la rédaction scientifique}
\addcontentsline{toc}{section}{2. Avancement de la rédaction scientifique}

Deux travaux de rédaction sont actuellement en cours. Leur avancement est détaillé
ci-dessous de façon transparente, y compris les difficultés rencontrées.

\subsection*{2a. Article issu du mémoire de Master}

Cet article porte sur le framework d'apprentissage fédéré centré sur les données
développé dans le cadre du mémoire de Master, appliqué à la maintenance prédictive
des éoliennes à partir de données SCADA. L'algorithme d'agrégation utilisé est
\textbf{FedProx}, choisi pour sa robustesse face à l'hétérogénéité des données
entre clients.

\textbf{État d'avancement par section :}

\begin{table}[h!]
\centering
\renewcommand{\arraystretch}{1.5}
\begin{tabularx}{\textwidth}{lXl}
\toprule
Section & Contenu & Statut \\
\midrule
Introduction & Contexte, motivation, contributions & \textcolor{orange}{\textbf{En attente}} \\
Related Work & Revue FL + maintenance éolienne & \textcolor{green!60!black}{\textbf{Rédigée}} \\
Méthodologie & Framework data-centric + FedProx & \textcolor{green!60!black}{\textbf{Rédigée}} \\
Expériences & Setup, datasets SCADA, protocoles & \textcolor{green!60!black}{\textbf{Rédigée}} \\
\textbf{Résultats} & \textbf{Classification multi-classe, évaluation FL} & \textcolor{gapred}{\textbf{Bloquée}} \\
Conclusion & Synthèse et perspectives & \textcolor{orange}{\textbf{En attente}} \\
\bottomrule
\end{tabularx}
\end{table}

\textbf{Description du problème de résultats :}

Deux difficultés majeures bloquent la finalisation de cette section :

\begin{enumerate}[leftmargin=*, label=\textbf{(\roman*)}]
  \item \textbf{Faible performance du modèle central.} Les modèles de classification
  multi-classe testés (sur les données SCADA éoliennes) n'atteignent pas des niveaux
  de précision satisfaisants. Cela peut être lié à la qualité du prétraitement des
  données, à un déséquilibre de classes non traité, ou à un choix de modèle inadapté.
  Des investigations sont en cours pour identifier la cause exacte.

  \item \textbf{Impossibilité d'obtenir des résultats d'entraînement fédéré.}
  Le pipeline d'entraînement FL (FedProx) ne produit pas encore de résultats
  stables ou exploitables. Le processus de refactorisation et de débogage est
  en cours.
\end{enumerate}

\begin{notebox}
\textbf{Plan d'action pour débloquer cet article :}
\begin{itemize}[noitemsep]
  \item Vérifier et corriger le pipeline de prétraitement des données SCADA
  (normalisation, gestion des valeurs manquantes, équilibrage des classes).
  \item Tester des modèles de référence simples (Random Forest, SVM) avant
  les modèles profonds pour établir un plancher de performance.
  \item Déboguer l'implémentation FedProx pas à pas, en commençant par un
  scénario à deux clients seulement.
  \item Cibler \textbf{IEEE Transactions on Industrial Informatics (TII)} ou
  \textbf{Renewable Energy} (Q1, Elsevier) pour la soumission finale.
\end{itemize}
\end{notebox}

\subsection*{2b. Revue systématique de littérature (SLR)}

Ce deuxième article vise à produire une revue systématique de la littérature sur
les \textit{Federated Time Series Foundation Models} pour la maintenance industrielle,
en suivant la méthodologie PRISMA (\textit{Preferred Reporting Items for Systematic
Reviews and Meta-Analyses}).

\textbf{Protocole PRISMA — Résultats préliminaires :}

La démarche PRISMA repose sur un processus d'entonnoir : identifier un grand nombre
de documents, puis les filtrer progressivement pour ne retenir que les plus pertinents.

\begin{table}[h!]
\centering
\caption{Résultats du processus de sélection PRISMA}
\label{tab:prisma-results}
\newcolumntype{R}{>{\raggedleft\arraybackslash}X}
\newcolumntype{C}{>{\centering\arraybackslash}X}
\renewcommand{\arraystretch}{1.5}
\begin{tabularx}{\textwidth}{lXC}
\toprule
Étape & Description & Résultat \\
\midrule
Identification & Requêtes sur IEEE Xplore, Scopus, Web of Science &
\textbf{3,000} \\
Dédoublonnage & Suppression des doublons entre bases &
\textbf{1,700} \\
Criblage des titres & Exclusion des articles hors sujet sur titre/résumé &
\textbf{361} \\
Éligibilité & Lecture complète, vérification des critères &
\textbf{181} \\
\textbf{Inclusion finale} & \textbf{Articles retenus pour la synthèse} &
\textbf{181} \\
\bottomrule
\end{tabularx}
\end{table}

\textbf{Requêtes utilisées} (exemples) :
\begin{itemize}[noitemsep]
  \item \textit{``Federated Learning'' AND ``Time Series'' AND ``Foundation Model''}
  \item \textit{``Federated'' AND ``Predictive Maintenance'' AND ``Deep Learning''}
  \item \textit{``Time Series Foundation Model'' AND ``Industrial''}
  \item \textit{``FedAvg'' OR ``FedProx'' AND ``SCADA'' OR ``vibration'' OR ``anomaly detection''}
\end{itemize}

\textbf{État actuel :} une liste de références pertinentes a été constituée mais n'est
pas encore formellement structurée selon la taxonomie PRISMA. La prochaine étape est
de formaliser cette liste et de compléter les chiffres du tableau ci-dessus.

% -------------------------------------------------------
\section*{3. Activités d'enseignement}
\addcontentsline{toc}{section}{3. Activités d'enseignement}

En parallèle de la recherche, des responsabilités pédagogiques ont été assumées
en tant que chargé de Travaux Dirigés (TD) et de Travaux Pratiques (TP) :

\begin{table}[h!]
\centering
\renewcommand{\arraystretch}{1.4}
\begin{tabularx}{\textwidth}{lXl}
\toprule
Module & Description & Type d'intervention \\
\midrule
Traitement du Signal & Encadrement des séances de TD pour le module de
traitement du signal pour la premier annee licence branch AI. Préparation des exercices, correction, et
accompagnement des étudiants. & TD \\[4pt]
Machine Learning & Encadrement des séances de TP pour un module d'apprentissage
automatique, pour les etudiants de la premier annee licence branch PC. Aide aux étudiants sur les implémentations pratiques,
correction des travaux et démonstrations en salle. & TP \\
\bottomrule
\end{tabularx}
\end{table}

Ces activités représentent un investissement en temps non négligeable. Elles
développent cependant des compétences de communication et de pédagogie qui
bénéficient directement à la qualité de la recherche, notamment pour la présentation
des travaux lors de conférences.

% ============================================================
% CHAPTER 1 — CONTEXTE GÉNÉRAL
% ============================================================
\chapter{Contexte général du sujet}
\markboth{Contexte général du sujet}{}

\section{L'évolution de la maintenance industrielle}

Pendant longtemps, les industries ont géré leurs équipements de deux façons simples :
soit on attend la panne pour réparer (\textit{maintenance corrective}), soit on
effectue des vérifications à intervalles réguliers, indépendamment de l'état réel
de la machine (\textit{maintenance préventive}). Ces deux approches ont leurs limites.
La première expose à des arrêts non planifiés coûteux ; la seconde gaspille du temps
et des ressources sur des équipements qui n'en ont pas besoin.

L'avènement des capteurs industriels et de l'apprentissage automatique a rendu possible
une troisième voie : la \textbf{maintenance prédictive}. L'idée est d'observer en
continu les signaux produits par les machines (vibrations, température, consommation
électrique, données SCADA) et d'en déduire, avant la panne, que quelque chose ne va pas.
Plutôt que de subir la défaillance ou d'intervenir inutilement, on agit au bon moment.

\begin{table}[h!]
\centering
\caption{Évolution des approches de maintenance}
\label{tab:maintenance-evolution}
\renewcommand{\arraystretch}{1.5}
\begin{tabularx}{\textwidth}{lXX}
\toprule
Approche & Principe & Limite principale \\
\midrule
Corrective & On répare après la panne & Arrêts non planifiés, coûts élevés \\
Préventive & On intervient à intervalles fixes & Gaspillage, interventions inutiles \\
Prédictive & On anticipe la panne via les données & Nécessite données et modèles IA \\
Prescriptive & L'IA recommande automatiquement l'action & Niveau de maturité encore faible \\
\bottomrule
\end{tabularx}
\end{table}

Les chiffres illustrent l'importance de ce changement de paradigme : on estime que
70\,\% des pannes industrielles sont précédées de signaux détectables dans les données,
et que la maintenance prédictive peut réduire les coûts de maintenance de 40\,\% dans
certains secteurs. Le marché mondial de cette technologie est estimé à plus de 630
milliards de dollars d'ici 2028.

\section{Le problème fondamental : des données utiles mais inaccessibles}

La maintenance prédictive dépend de la quantité et de la qualité des données
disponibles pour entraîner les modèles. En théorie, plus on centralise des données
provenant de nombreux sites et équipements, meilleurs sont les modèles. En pratique,
cette centralisation est rarement possible, pour plusieurs raisons :

\begin{itemize}
  \item \textbf{Confidentialité et concurrence.} Les données de production ou de
  performance d'une usine sont stratégiquement sensibles. Aucun opérateur industriel
  ne souhaite partager ses données brutes avec ses concurrents.
  \item \textbf{Réglementations.} Le RGPD et diverses réglementations sectorielles
  encadrent strictement les transferts de données, notamment dans l'énergie ou la défense.
  \item \textbf{Volume.} Un parc d'éoliennes peut générer des dizaines de gigaoctets
  de données par jour. Centraliser tout cela n'est ni pratique, ni économique.
  \item \textbf{Hétérogénéité.} Chaque domaine industriel produit des données différentes :
  un signal SCADA d'éolienne n'a rien à voir avec les données acoustiques d'une ligne
  de production ou avec les mesures thermiques d'un système HVAC.
\end{itemize}

Cette situation crée un paradoxe : les données les plus utiles pour améliorer la
maintenance sont précisément celles qu'on ne peut pas partager. La question centrale
devient alors : \textit{comment entraîner de bons modèles sans jamais voir les données
brutes des autres ?}

\section{Positionnement de la thèse}

C'est exactement à cette question que cette thèse cherche à répondre, en combinant
trois champs de recherche :

\begin{itemize}
  \item \textbf{Le Federated Learning (FL)}, qui permet d'entraîner des modèles
  collectivement sans centraliser les données.
  \item \textbf{Les Foundation Models pour séries temporelles}, qui permettent de
  construire des modèles puissants, pré-entraînés sur de grandes quantités de données,
  adaptables à de nombreuses tâches avec peu d'effort.
  \item \textbf{La maintenance prédictive industrielle multi-domaine}, qui est
  l'application cible : anticiper les pannes sur des équipements variés (éoliennes,
  systèmes HVAC, machines de production) sans partager les données de terrain.
\end{itemize}

\begin{infobox}
\textbf{La proposition centrale de la thèse :} développer un framework basé sur des
Foundation Models de séries temporelles, entraînés de façon fédérée sur des données
industrielles hétérogènes, capable de s'adapter à plusieurs domaines avec un minimum
de personnalisation --- et ce, sans que les données brutes ne quittent jamais les sites locaux.
\end{infobox}

Cette vision n'est pas encore réalisée dans la littérature. C'est précisément l'espace
blanc que cette thèse vise à combler.

% ============================================================
% CHAPTER 2 — CONCEPTS DE BASE
% ============================================================
\chapter{Concepts de base}
\markboth{Concepts de base}{}

Ce chapitre présente les quatre piliers théoriques qui structurent le sujet de thèse.
L'objectif n'est pas de les définir de façon exhaustive, mais de donner à tout lecteur
une compréhension claire et fonctionnelle de chacun.

\section{L'apprentissage fédéré (Federated Learning)}

L'apprentissage fédéré est une méthode d'entraînement de modèles de machine learning
dans laquelle les données restent sur les appareils ou sites qui les produisent. Au lieu
d'envoyer les données vers un serveur central, on envoie uniquement les \textit{mises à
jour du modèle} (les poids, les gradients).

Le déroulement classique est le suivant :
\begin{enumerate}
  \item Un serveur central distribue un modèle global à tous les clients.
  \item Chaque client entraîne ce modèle localement sur ses propres données.
  \item Chaque client envoie au serveur uniquement les mises à jour calculées,
  pas les données brutes.
  \item Le serveur agrège ces mises à jour (par exemple en faisant une moyenne
  pondérée) pour produire un nouveau modèle global amélioré.
  \item Le cycle recommence.
\end{enumerate}

\textbf{FedProx} est l'algorithme utilisé dans le travail issu du Master. C'est une
extension de l'algorithme standard (FedAvg) qui ajoute un mécanisme de contrôle pour
éviter que les clients ne s'éloignent trop du modèle global lors de l'entraînement local.
Cela le rend plus robuste quand les données des différents clients sont très dissemblables
(situation appelée \textit{non-IID}).

Le défi majeur en FL est précisément cette hétérogénéité : les données de chaque client
peuvent avoir des distributions, des formats ou des volumes très différents. C'est
particulièrement vrai dans le contexte industriel, où chaque site a ses propres capteurs,
ses propres conditions d'exploitation et ses propres types de pannes.

\section{Les séries temporelles industrielles}

Une série temporelle est une séquence de mesures enregistrées à intervalles réguliers
dans le temps. Dans le contexte industriel, cela couvre les signaux de capteurs de
vibration, de température, de pression, de consommation électrique, ou encore les
données SCADA des éoliennes.

Ces séries présentent des caractéristiques qui compliquent leur analyse :

\begin{itemize}
  \item Elles sont souvent \textbf{multivariées} (plusieurs capteurs simultanément).
  \item Elles sont \textbf{non-stationnaires} : les conditions changent dans le temps
  (charge de la machine, saison, usure).
  \item Les défaillances sont des \textbf{événements rares} : les données sont
  fortement déséquilibrées (peu de pannes, beaucoup d'état normal).
  \item Elles contiennent du \textbf{bruit}, des valeurs manquantes et des anomalies
  de capteurs qui ne correspondent pas à de vraies pannes.
\end{itemize}

Dans un cadre fédéré, ces difficultés sont amplifiées : chaque client peut avoir
des fréquences d'échantillonnage différentes, des longueurs de séries variables,
et des types de défaillances propres à ses équipements.

\section{La maintenance prédictive}

La maintenance prédictive consiste à utiliser les données issues des capteurs pour
anticiper l'état futur d'un équipement et planifier les interventions au bon moment.
Les principales tâches associées sont :

\begin{itemize}
  \item \textbf{Détection d'anomalies} : identifier un comportement anormal.
  \item \textbf{Classification de défaillances} : reconnaître le type de panne
  parmi plusieurs catégories possibles (tâche principale de cette thèse).
  \item \textbf{Estimation de la durée de vie résiduelle (RUL)} : prédire combien
  de temps il reste avant la panne.
  \item \textbf{Prévision de l'état de santé} : suivre l'évolution globale d'un
  indicateur de santé machine.
\end{itemize}

\section{Les Foundation Models pour séries temporelles}

Un Foundation Model est un grand modèle pré-entraîné sur un corpus massif et
hétérogène de données, puis adapté à diverses tâches en aval avec peu d'exemples
et peu de réentraînement. Le principe vient du traitement du langage naturel :
GPT ou BERT sont des Foundation Models pour le texte. L'idée a depuis été étendue
aux images, puis aux séries temporelles.

Les principaux Foundation Models pour séries temporelles apparus entre 2023 et 2024 sont :

\begin{table}[h!]
\centering
\caption{Principaux Foundation Models pour séries temporelles (2023--2024)}
\label{tab:tsfm-comparison}
\renewcommand{\arraystretch}{1.5}
\begin{tabularx}{\textwidth}{llXl}
\toprule
Modèle & Année & Architecture et particularités & Tâches principales \\
\midrule
TimesFM & 2024 & Transformer décodeur, pré-entraîné sur 100\,Mds de points
& Prévision \\
Moirai & 2024 & Encodeur masqué avec mélange d'experts, multi-fréquence
& Prévision + anomalie \\
MOMENT & 2024 & Encodeur-décodeur (type T5), famille de modèles pour tâches multiples
& Prévision + classification \\
Lag-Llama & 2023 & Transformer décodeur, probabiliste, univarié
& Prévision \\
UniTime & 2023 & Transformer avec prompts de domaine, multi-domaine
& Prévision \\
TimeLLM & 2022 & LLM (GPT-2) réutilisé pour les séries temporelles
& Prévision \\
\bottomrule
\end{tabularx}
\end{table}

Ce qui est révolutionnaire avec ces modèles, c'est leur capacité à fonctionner
en \textbf{zero-shot} ou \textbf{few-shot} : sans entraînement spécifique sur un
nouveau domaine, ils peuvent déjà produire des prédictions raisonnables. Cela ouvre
la possibilité d'un modèle unique capable de servir de nombreux équipements industriels
différents.

\begin{infobox}
\textbf{Le chaînon manquant :} aucun de ces Foundation Models n'a été conçu pour être
entraîné de façon fédérée. Et aucun d'eux n'a été évalué sur des données industrielles
réelles (vibration, SCADA, acoustique). C'est précisément le double gap que cette thèse explore.
\end{infobox}

% ============================================================
% CHAPTER 3 — ÉTAT DE L'ART
% ============================================================
\chapter{État de l'art}
\markboth{État de l'art}{}

\section{Cartographie du domaine de recherche}

La recherche sur l'apprentissage fédéré de Foundation Models pour séries temporelles
s'articule autour de trois axes principaux, illustrés dans la figure~\ref{fig:landscape}.

\begin{figure}[h!]
\centering
\begin{tikzpicture}[
    node distance=1.2cm,
    box/.style={rectangle, rounded corners, draw=mainblue, fill=mainblue!10,
                text width=3.2cm, minimum height=1cm, align=center, font=\small},
    central/.style={rectangle, rounded corners, draw=gapred, fill=gapred!15,
                   text width=4cm, minimum height=1.3cm, align=center, font=\bfseries},
    arrow/.style={->, >=stealth, thick, mainblue}
]
% Central node
\node[central] (center) at (0,0) {FL + TS FM +\\Maintenance};

% Surrounding nodes
\node[box, above=1.5cm of center] (tsfm) {Foundation Models\\Séries Temporelles};
\node[box, left=1.8cm of center] (fl) {Apprentissage\\Fédéré};
\node[box, right=1.8cm of center] (pm) {Maintenance\\Prédictive};
\node[box, below left=1.2cm and 1.5cm of center] (nlp) {Techniques NLP/Vision\\(LoRA, Adapters)};
\node[box, below right=1.2cm and 1.5cm of center] (bench) {Benchmarks\\Évaluation};

% Arrows
\draw[arrow] (tsfm) -- (center);
\draw[arrow] (fl) -- (center);
\draw[arrow] (pm) -- (center);
\draw[arrow, dashed, gaporange] (nlp) -- (center);
\draw[arrow, dashed, gaporange] (bench) -- (center);

% Maturity labels
\node[font=\tiny, above=0.1cm of tsfm] {\textcolor{gaporange}{Émergent (~8 modèles)}};
\node[font=\tiny, left=0.1cm of fl] {\textcolor{mainblue}{Mature}};
\node[font=\tiny, right=0.1cm of pm] {\textcolor{mainblue}{Classique}};
\node[font=\tiny, below=0.1cm of center] {\textcolor{gapred}{Vierge (3 travaux)}};

\end{tikzpicture}
\caption{Cartographie des domaines de recherche et leur intersection. Les flèches pleines
indiquent les contributions directes, les flèches pointillées les inspirations méthodologiques.}
\label{fig:landscape}
\end{figure}

\begin{table}[h!]
\centering
\caption{Niveaux de maturité de la littérature par axe de recherche}
\label{tab:maturity}
\renewcommand{\arraystretch}{1.6}
\begin{tabularx}{\textwidth}{>{\centering\arraybackslash}p{2.5cm}X>{\centering\arraybackslash}p{2cm}>{\centering\arraybackslash}p{2.5cm}}
\toprule
\textbf{Domaine} & \textbf{Description} & \textbf{Maturité} & \textbf{Nb. travaux} \\
\midrule
\rowcolor{mainblue!10}
FL général & Algorithmes d'agrégation (FedAvg, FedProx), hétérogénéité, privacy & Mature & 1000+ \\
\rowcolor{gaporange!10}
TS FMs centr. & TimesFM, Moirai, MOMENT, Lag-Llama, UniTime & Émergent & ~8 modèles \\
\rowcolor{gaporange!10}
FL pour NLP/Vision & LoRA fédéré, adapters, split learning, quantization & Émergent & ~50 travaux \\
\rowcolor{gapred!10}
FL pour maintenance & Pruckovskaja et al. (2023), OASEES (2026), turbofan 2024 & Initial & ~10 travaux \\
\rowcolor{gapred!20}
\textbf{FL + TS FM} & \textbf{Time-FFM, Ali et al., FedForecaster} & \textbf{Quasi vide} & \textbf{3 travaux} \\
\rowcolor{gapred!25}
\textbf{FL + TS FM + PM} & \textbf{Aucun travail publié} & \textbf{Vierge} & \textbf{0} \\
\bottomrule
\end{tabularx}
\end{table}

\section{Foundation Models pour séries temporelles}

\subsection{Panorama des modèles existants (2023--2024)}

Les Foundation Models pour séries temporelles constituent un domaine émergent, avec
la publication entre 2023 et 2024 de plusieurs architectures majeures pré-entraînées
sur de larges corpus de données publiques.

\begin{table}[h!]
\centering
\caption{Comparaison des principaux Foundation Models pour séries temporelles}
\label{tab:tsfm-detailed}
\renewcommand{\arraystretch}{1.5}
\begin{tabularx}{\textwidth}{llXp{2.5cm}l}
\toprule
\textbf{Modèle} & \textbf{Année} & \textbf{Architecture et particularités} & \textbf{Tâches} & \textbf{Params} \\
\midrule
TimesFM & 2024 & Transformer décodeur, pré-entraîné sur 100Mds points, contexte 16k & Prévision & 200M \\
Moirai & 2024 & Encodeur masqué, MoE, multi-fréquence, LOTSA (320M pts) & Prév. + Anomalie & 32--480M \\
MOMENT & 2024 & Architecture T5-like, Time-Series Pile, multi-tâches & Prév. + Classif. & 385M \\
Lag-Llama & 2023 & Décodeur probabiliste, features lagged, univarié & Prévision & ~100M \\
UniTime & 2023 & Domain prompts, multi-domaine, Language-TS Transformer & Prévision & ND \\
TimeLLM & 2022 & GPT-2 reprogrammé, tokenization texte, LoRA fine-tuning & Prévision & 345M \\
Timer-S1 & 2026 & Sparse MoE (8.3B total), TimeBench (1T pts), contexte 11.5k & Prévision & 750M actifs \\
\bottomrule
\end{tabularx}
\end{table}

\subsection{Limitations dans le contexte industriel}

Ces modèles présentent quatre limitations majeures pour la maintenance prédictive :

\begin{enumerate}
  \item \textbf{Distribution des données d'entraînement.} Pré-entraînement sur données
  publiques (météo, trafic web, finance) via LOTSA ou TimeBench, sans données industrielles
  réelles (vibrations, SCADA, acoustique).

  \item \textbf{Support multivarié limité.} Conçus principalement pour séries univariées
  ou faible dimension, alors que les données industrielles comportent souvent 100+ capteurs.

  \item \textbf{Absence d'évaluation maintenance.} Aucun test sur classification de
  défaillances ou détection d'anomalies industrielles.

  \item \textbf{Non-conception pour le FL.} Aucune propriété \textit{federation-friendly}
  (patching adaptatif, mise à jour sparse, quantization-aware).
\end{enumerate}

\begin{infobox}
\textbf{Écart de performance estimé :} Les tests préliminaires suggèrent une perte de
performance de 25--50\% lors de l'application de TS FMs sur données industrielles réelles
sans fine-tuning spécifique.
\end{infobox}

\section{Federated Learning pour Foundation Models}

\subsection{Stratégies de personnalisation fédérée (NLP/Vision)}

La communauté NLP/Vision a développé des techniques d'adaptation efficace en paramètres
(PEFT) pour le FL de grands modèles. Le tableau suivant synthétise les approches
potentiellement transférables aux séries temporelles.

\begin{table}[h!]
\centering
\caption{Stratégies de personnalisation fédérée pour Foundation Models (inspirations NLP/Vision)}
\label{tab:peft-strategies}
\renewcommand{\arraystretch}{1.6}
\begin{tabularx}{\textwidth}{lp{3cm}Xl}
\toprule
\textbf{Méthode} & \textbf{Auteurs (Année)} & \textbf{Principe} & \textbf{Gain} \\
\midrule
FedCLIP & Wang et al. (2023) & Adapters sur backbone gelé, agrégation des adapters uniquement & 283$\times$ vitesse \\
FedOPAL & Tupper \& Gagné (2025) & Dual LoRA (global + personnel) avec orthogonalité & Personnalisation \\
FDLoRA & Qi et al. (2024) & Fusion adaptative global/personnel post-agrégation & Meilleure perf. \\
HeLoRA & Fan et al. (2025) & Rangs LoRA hétérogènes selon capacités clients & Flexibilité \\
FedQLoRA & Hu et al. (2025) & Adaptateur quantization-aware & Robustesse \\
FedsLLM & Zhao et al. (2024) & Split learning + LoRA, optimisation delay & 47.6\% réduction delay \\
\bottomrule
\end{tabularx}
\end{table}

\subsection{Transposition aux séries temporelles}

L'adaptation de ces techniques aux TS FMs soulève des défis spécifiques :

\begin{itemize}
  \item \textbf{Adapters temporels :} Les adapters doivent capturer des patterns
  temporels (attention sur le temps) plutôt que des features textuelles.
  \item \textbf{Domain prompts :} UniTime et Time-FFM montrent l'intérêt des prompts
  de domaine pour l'adaptation cross-domaine en séries temporelles.
  \item \textbf{Patching fédéré :} Le mécanisme de patching (découpage en patches)
  pourrait être adapté par client selon sa fréquence d'échantillonnage locale.
\end{itemize}

\section{Federated Time Series Foundation Models : les trois travaux existants}

\subsection{Tableau comparatif des travaux à l'intersection}

\begin{table}[h!]
\centering
\caption{Comparaison des travaux FL + TS Foundation Models}
\label{tab:fl-tsfm-works}
\renewcommand{\arraystretch}{1.6}
\begin{tabularx}{\textwidth}{>{\bfseries}p{2.5cm}XXXX}
\toprule
\textbf{Travail} & \textbf{Time-FFM} & \textbf{Ali et al.} & \textbf{FedForecaster} \\
\midrule
Auteurs & Liu et al. (2024) & Ali et al. (2025) & Maher et al. (2025) \\
Venue & NeurIPS & arXiv & EDBT \\
Approche & FM fédéré natif & Fine-tuning fédéré & AutoML fédéré \\
Modèles & LLM reprogrammé & Chronos, TimesFM & Méta-modèle \\
Domaine & Générique & Médical (ECG) & Générique \\
Personnalisation & Têtes locales & FedAvg standard & Recommandation \\
Maintenance & Non & Non & Non \\
\bottomrule
\end{tabularx}
\end{table}

\subsection{Analyse critique par travail}

\textbf{Time-FFM} représente la première tentative de Foundation Model fédéré pour
séries temporelles. L'approche reprogramme un LLM pré-entraîné via tokenization texte
des séries, avec des prompts de domaine dynamiques et des têtes de prédiction locales
par client. Les résultats montrent des capacités zero-shot et few-shot supérieures aux
baselines. \textbf{Limites :} données génériques uniquement, pas de tâches de
classification (maintenance), pas de techniques PEFT (LoRA).

\textbf{Ali et al.} étudient le fine-tuning fédéré de modèles pré-entraînés (Chronos,
TimesFM) sur données ECG médicales. L'étude montre que FL peut dépasser le fine-tuning
local sous hétérogénéité modérée, mais les gains disparaissent sous forte non-IID.
\textbf{Limites :} domaine médical uniquement, pas de solution à l'hétérogénéité,
algorithme FedAvg standard.

\textbf{FedForecaster} propose un système AutoML fédéré où les clients envoient des
méta-caractéristiques statistiques, et un serveur apprend un méta-modèle pour
recommander les algorithmes de prévision. Cette approche est légère mais ne permet
pas de transfert de connaissance cross-domaine via un FM monolithique. \textbf{Limites :}
pas un vrai FM, pas de maintenance industrielle.

\section{Federated Learning et Maintenance Prédictive}

\subsection{Travaux FL appliqués à la maintenance}

L'application du FL à la maintenance prédictive reste classique (sans Foundation Models) :

\begin{table}[h!]
\centering
\caption{Travaux de FL pour maintenance prédictive}
\label{tab:fl-pm-works}
\renewcommand{\arraystretch}{1.5}
\begin{tabularx}{\textwidth}{p{3.5cm}p{2.5cm}Xl}
\toprule
\textbf{Travail} & \textbf{Application} & \textbf{Contribution} & \textbf{Année} \\
\midrule
Pruckovskaja et al. & Pompes, moteurs & Évaluation FL sur 4 datasets industriels, performance data-dépendante & 2023 \\
OASEES SWARMAware & Éoliennes (acoustique) & Détection défauts sur capteur, généralisation fleet-wide & 2026 \\
Turbofan RUL FL & Moteurs turbofan & RUL prediction via FATE/DC-Fed, performance comparable centralisé & 2024 \\
\bottomrule
\end{tabularx}
\end{table}

\subsection{Constat majeur : absence de Foundation Models}

Aucun des travaux de maintenance prédictive en FL n'utilise de Foundation Model. La
combinaison des trois éléments --- \textbf{FL + TS Foundation Model + Maintenance
Prédictive Industrielle} --- constitue un espace de recherche entièrement vierge.

% ============================================================
% CHAPTER 4 — PROBLÉMATIQUES IDENTIFIÉES
% ============================================================
\chapter{Problématiques identifiées}
\markboth{Problématiques identifiées}{}

L'analyse de la littérature fait émerger quatre problématiques fondamentales. Elles
ne sont pas indépendantes : chacune bloque l'autre, et c'est leur combinaison qui
explique pourquoi le sujet est encore vierge.

\section{P1 — La personnalisation fédérée sur données hétérogènes}

\textbf{Le problème.} Les clients d'un système FL industriel ont des données très
différentes les unes des autres : fréquences d'échantillonnage distinctes, types de
capteurs variés, saisonnalités propres à chaque site. C'est ce qu'on appelle le
problème \textit{non-IID}. La grande taille des Foundation Models amplifie encore
ce phénomène, car les mises à jour locales peuvent diverger fortement du modèle global.
Les stratégies de personnalisation classiques en FL (têtes de prédiction locales)
ne capturent pas les dépendances temporelles profondes propres à chaque client.

\textbf{Ce qui existe.} Time-FFM (2024) propose une tête locale par client plus un
encodeur global. LoRA-TS (2024) applique LoRA aux séries temporelles en setting
centralisé uniquement. Aucun travail ne combine prompts de domaine, LoRA et FL pour
les séries temporelles.

\textbf{Ce qui manque.} Un mécanisme de personnalisation fédérée adapté aux séquences
temporelles, combinant des prompts spécifiques au domaine client et du fine-tuning
efficace en paramètres (LoRA sur couches d'attention temporelle).

\begin{infobox}
\textbf{Question de recherche RQ1 :} Comment concevoir un mécanisme de personnalisation
fédérée, combinant prompts domaine-spécifiques et LoRA, permettant à un Time Series
Foundation Model de s'adapter aux distributions locales hétérogènes de chaque client
industriel ?
\end{infobox}

\section{P2 — L'absence de benchmark fédéré standardisé}

\textbf{Le problème.} Il est impossible de comparer ou de valider des contributions
de recherche dans ce domaine s'il n'existe pas de benchmark commun. FoundTS (2024)
est le seul benchmark pour les TS FMs, mais il est entièrement centralisé, focalisé
sur la prévision, et ne propose aucune métrique adaptée à la maintenance (F1 sur classes
déséquilibrées, AUC-PR). Sans protocole standardisé, chaque travail est une île.

\textbf{Ce qui existe.} FoundTS (2024) pour le centralisé. Des métriques d'anomalie
sur les datasets UCR/ICS. Aucun protocole fédéré multi-sites pour la maintenance.

\textbf{Ce qui manque.} Un benchmark fédéré avec partitions non-IID réalistes,
métriques de dérive client, et évaluation adaptée aux tâches de maintenance.

\begin{infobox}
\textbf{Question de recherche RQ2 :} Quels protocoles d'évaluation et quels benchmarks
fédérés permettent de mesurer fidèlement les capacités de généralisation cross-domaine
d'un TS FM dans un contexte de maintenance prédictive industrielle multi-sites ?
\end{infobox}

\section{P3 — L'incompatibilité architecturale des TS FMs avec le FL}

\textbf{Le problème.} Les architectures actuelles des TS FMs n'ont pas été conçues
avec les contraintes FL en tête. Les Transformers décodeurs (comme TimesFM) requièrent
des contextes longs (plus de 10\,000 tokens), ce qui est prohibitif en communication
FL. Les mécanismes de mélange d'experts (MoE) divergent sous données non-IID.
L'attention causale résiste à l'agrégation globale.

\textbf{Ce qui existe.} Le patching (Moirai, MOMENT) réduit la longueur de contexte
et est plus compatible FL. Aucun TS FM n'a été conçu nativement pour la fédération.

\textbf{Ce qui manque.} Une architecture intégrant des propriétés \textit{federation-friendly}
dès sa conception : patching adaptatif, MoE sparse guidé par le client, et mises à
jour low-rank limitées aux couches clés.

\begin{infobox}
\textbf{Question de recherche RQ3 :} Quelles propriétés architecturales d'un Time Series
Foundation Model — tokenisation, mécanisme d'attention, stratégie de mise à jour — maximisent
son efficacité de communication et sa convergence dans un environnement FL hétérogène ?
\end{infobox}

\section{P4 — L'inadéquation aux données industrielles réelles}

\textbf{Le problème.} Les TS FMs sont pré-entraînés sur des données publiques génériques
(météo, finance, trafic web — corpus LOTSA). Les signaux industriels ont des
caractéristiques radicalement différentes : haute fréquence, non-stationnarité forte,
bruit de capteurs, contraintes physiques. Ni TimesFM, ni Moirai, ni MOMENT n'ont
jamais été testés sur des données SCADA, vibration ou acoustique industrielle.
L'écart de performance estimé sur des données réelles est de 25 à 50\,\%.

\textbf{Ce qui existe.} Des tests sur données électriques (proxy trop faible).
Pas d'évaluation sur données de maintenance réelle.

\textbf{Ce qui manque.} Un corpus de pré-entraînement industriel et une évaluation
zero-shot et few-shot sur des tâches de maintenance réelles.

\begin{infobox}
\textbf{Question de recherche RQ4 :} Dans quelle mesure les Time Series Foundation Models
actuels peuvent-ils se généraliser aux signaux industriels (SCADA, vibration, acoustique)
en zero-shot ou few-shot, et comment l'entraînement fédéré multi-domaine peut-il réduire
cet écart ?
\end{infobox}

\section{Synthèse : carte des gaps}

\begin{table}[h!]
\centering
\caption{Synthèse des gaps identifiés}
\label{tab:gaps}
\renewcommand{\arraystretch}{1.6}
\begin{tabularx}{\textwidth}{>{\raggedright\arraybackslash}p{4.5cm}X>{\centering\arraybackslash}p{2.2cm}>{\raggedright\arraybackslash}p{4cm}}
\toprule
\textbf{Gap} & \textbf{Description} & \textbf{Sévérité} & \textbf{Partiellement adressé ?} \\
\midrule
Entraînement fédéré de TS FMs & Aucune méthode complète & \textcolor{gapred}{\textbf{Critique}} & Time-FFM (prévision seulement) \\
TS FMs pour maintenance & Aucun travail existant & \textcolor{gapred}{\textbf{Critique}} & Non \\
TS FMs sur non-IID hétérogène & Chute de perf. sous forte hétéro. & \textcolor{gapred}{\textbf{Critique}} & Ali et al. (partiel, médical) \\
PEFT (LoRA/adapters) pour TS FMs & Inexploré en FL & \textcolor{gapred}{\textbf{Critique}} & Non \\
Généralisation cross-domaine fédérée & Seulement Vision/NLP & \textcolor{gaporange}{\textbf{Majeur}} & FedAG (vision uniquement) \\
Benchmark fédéré multi-domaine & Aucun & \textcolor{gaporange}{\textbf{Majeur}} & Non \\
Privacy (DP + FL + TS FMs) & Inexploré & \textcolor{gapyellow}{\textbf{Modéré}} & Non \\
\bottomrule
\end{tabularx}
\end{table}

% ============================================================
% CHAPTER 5 — CONTRIBUTIONS ENVISAGÉES
% ============================================================
\chapter{Contributions envisagées}
\markboth{Contributions envisagées}{}

Les quatre questions de recherche identifiées au chapitre précédent définissent
naturellement quatre contributions complémentaires, couvrant l'ensemble du doctorat
et progressant du plus fondamental (méthode) au plus concret (validation industrielle).

\section{Vue d'ensemble du plan de contributions}

\begin{table}[h!]
\centering
\caption{Plan des contributions sur 3 ans}
\label{tab:contributions}
\renewcommand{\arraystretch}{1.5}
\begin{tabularx}{\textwidth}{llXll}
\toprule
\# & Nom & Description & RQ & Période \\
\midrule
C1 & FedPEFT-TS & Personnalisation fédérée par LoRA + prompts & RQ1 & An 1--2 \\
C2 & FedTS-Bench & Benchmark fédéré pour maintenance & RQ2 & An 1--2 \\
C3 & FedTSFM-Arch & Architecture TS FM native FL & RQ3 & An 2--3 \\
C4 & IndusFedFM & Validation industrielle multi-domaines & RQ4 & An 3 \\
\bottomrule
\end{tabularx}
\end{table}

Le fil conducteur est le suivant : partir d'un TS FM existant, le rendre personnalisable
en FL (C1), le valider sur un benchmark dédié (C2), proposer une architecture optimisée
(C3), puis valider le tout sur des données industrielles réelles multi-domaines (C4).

\section{C1 — FedPEFT-TS : personnalisation fédérée pour TS FMs}

\textbf{Problème adressé :} RQ1 — comment personnaliser un TS FM global pour des
clients industriels aux données très hétérogènes.

\textbf{Principe :}
\begin{itemize}
  \item On part d'un TS FM existant (Moirai ou TimesFM) et on gèle son backbone.
  \item On attache de petits modules LoRA sur les couches d'attention temporelle.
  \item Chaque client entraîne un LoRA \textit{personnel} (qui reste local) et
  un LoRA \textit{global} (qui est agrégé sur le serveur).
  \item Des prompts spécifiques au domaine client encodent le contexte de chaque site
  (inspiration : UniTime, Time-FFM).
\end{itemize}

\textbf{Ce que ça apporte par rapport à l'état de l'art :} Time-FFM (2024) utilise
des têtes locales mais pas de LoRA. LoRA-TS (2024) fonctionne en centralisé seulement.
FedPEFT-TS serait la première combinaison LoRA + prompts domaine + FL pour TS FMs.

\section{C2 — FedTS-Bench : benchmark fédéré pour la maintenance industrielle}

\textbf{Problème adressé :} RQ2 — absence totale de benchmark standardisé.

\textbf{Principe :}
\begin{itemize}
  \item Construire un protocole d'évaluation open-source à partir de datasets
  industriels publics : NASA C-MAPSS (moteurs turbofan), CWRU (roulements),
  FEMTO-ST (dégradation), données SCADA éoliennes.
  \item Définir des partitions non-IID réalistes simulant N clients industriels hétérogènes.
  \item Proposer des métriques adaptées : F1 pondéré, AUC-PR (classes déséquilibrées),
  variance de dérive client, nombre de rounds à convergence.
\end{itemize}

\textbf{Livrable concret :} splits de datasets + code d'évaluation disponible publiquement.
Ce benchmark serait la première ressource communautaire de ce type.

\section{C3 — FedTSFM-Arch : architecture TS FM native FL}

\textbf{Problème adressé :} RQ3 — aucun TS FM n'a été conçu avec les contraintes FL
en tête.

\textbf{Principe :} concevoir ou adapter un TS FM avec trois propriétés \textit{federation-friendly} :

\begin{table}[h!]
\centering
\renewcommand{\arraystretch}{1.5}
\begin{tabularx}{\textwidth}{lXX}
\toprule
Propriété & Mécanisme & Bénéfice FL \\
\midrule
Patching hiérarchique & Découpage adaptatif selon la fréquence locale du client &
Réduit la longueur de contexte, diminue la communication \\
MoE sparse client-guidé & Experts activés différemment selon le domaine client &
Spécialisation locale sans partage de données \\
Mises à jour low-rank & LoRA sur query/key uniquement &
Réduction potentielle $\times$50 du volume échangé \\
\bottomrule
\end{tabularx}
\end{table}

\section{C4 — IndusFedFM : validation industrielle multi-domaines}

\textbf{Problème adressé :} RQ4 — inadéquation des TS FMs aux données industrielles.

\textbf{Principe :} déployer le framework complet (C1 + C3) sur des cas d'usage
industriels réels multi-domaines (éolien SCADA, HVAC, manufacture) et évaluer
trois scénarios :

\begin{itemize}
  \item \textbf{Intra-domaine :} N clients éoliens hétérogènes (données SCADA du Master).
  \item \textbf{Cross-domaine :} éolien + HVAC + manufacture en FL simultané.
  \item \textbf{Cold-start :} nouveau client sans historique (test zero-shot).
\end{itemize}

\begin{infobox}
\textbf{Vision d'ensemble — FedTSFM :}
\begin{center}
C1 (méthode) + C2 (benchmark) + C3 (architecture) + C4 (validation)\\[4pt]
$\Downarrow$\\[4pt]
\textit{Un Foundation Model de Séries Temporelles entraînable de façon fédérée\\
pour la maintenance prédictive multi-domaine des systèmes industriels}
\end{center}
\end{infobox}

% ============================================================
% CHAPTER 6 — PUBLICATIONS PRÉVUES
% ============================================================
\chapter{Contenu des publications prévues}
\markboth{Publications prévues}{}

\section{Stratégie de publication}

La thèse vise cinq publications sur trois ans : une conférence pour le travail issu
du Master (P0), deux conférences ou journaux en première moitié de thèse (P1, P2),
et deux journaux Q1 pour les contributions majeures (P3, P4).

\begin{table}[h!]
\centering
\caption{Plan de publication complet}
\label{tab:publication-plan}
\renewcommand{\arraystretch}{1.6}
\begin{tabularx}{\textwidth}{>{\centering\arraybackslash}p{0.8cm}X>{\centering\arraybackslash}p{2.8cm}>{\centering\arraybackslash}p{2.2cm}>{\centering\arraybackslash}p{1.8cm}}
\toprule
\textbf{\#} & \textbf{Titre (provisoire)} & \textbf{Contribution} & \textbf{Type} & \textbf{Année} \\
\midrule
P0 & Data-Centric FL for Wind Turbine PdM & Master & Journal & 2025 \\
P1 & Literature Review for FL+TS FMs for Industrial Maintenance & SLR (C1--C2) & Journal & An 1--2 \\
P2 & FedTS-Bench: A Federated Benchmark for Industrial PM & C2 & Conférence & An 2 \\
P3 & FedTSFM-Arch: A Federation-Native TS FM Architecture & C3 & Journal Q1 & An 2--3 \\
P4 & IndusFedFM: Multi-Domain Federated TS FMs for PM & C4 & Journal Q1 & An 3 \\
\bottomrule
\end{tabularx}
\end{table}

\section{Détail par publication}

\textbf{P0 — Data-Centric Federated Learning for Wind Turbines Predictive Maintenance} (Journal, 2025)

L'article en cours porte sur le framework FL data-centric pour la maintenance prédictive
des éoliennes à partir de données SCADA, avec FedProx comme algorithme d'agrégation.
Les cibles de soumission sont IEEE Transactions on Industrial Informatics (TII) ou
Renewable Energy (Q1, Elsevier).

\textbf{P1 — Literature Review for Federated Time Series Foundation Models for Industrial Maintenance} (Journal, An 1--2)

Cette revue systématique de littérature (SLR) suivant la méthodologie PRISMA couvre
l'intersection entre l'apprentissage fédéré, les Foundation Models pour séries temporelles
et la maintenance prédictive industrielle. L'objectif est d'identifier les gaps de recherche
et de positionner les contributions de la thèse.

\textbf{P2 — FedTS-Bench} (Conférence, An 2)

L'article décrira la construction du benchmark (contribution C2), évaluera cinq
méthodes FL ou plus croisées avec trois TS FMs sur toutes les partitions, et analysera
la sensibilité au degré de non-IID. Le code et les splits seront rendus publics.

\textbf{P3 — FedTSFM-Arch} (Journal Q1, An 2--3)

L'article proposera la nouvelle architecture (contribution C3), avec une étude
d'ablation composant par composant et une comparaison contre TimesFM, Moirai et
MOMENT sur FedTS-Bench. L'analyse portera sur le compromis communication/performance.

\textbf{P4 — IndusFedFM} (Journal Q1, An 3)

L'article validera le framework complet sur données industrielles réelles multi-domaines
(contribution C4), avec les trois scénarios décrits (intra-domaine, cross-domaine,
cold-start). Il établira des baselines de référence pour la communauté et proposera
des recommandations pratiques pour le déploiement.

\section{Fil narratif de la thèse}

\begin{figure}[h!]
\centering
\begin{tikzpicture}[
    node distance=0.8cm,
    pubnode/.style={rectangle, rounded corners=3pt, draw=mainblue, fill=mainblue!8,
                    text width=12cm, minimum height=1.1cm, align=left,
                    font=\small, inner sep=6pt},
    yearnode/.style={rectangle, rounded corners=2pt, fill=mainblue!20,
                     text width=1.8cm, minimum height=0.6cm, align=center,
                     font=\tiny\bfseries, inner sep=2pt},
    arrow/.style={->, >=stealth, thick, mainblue, shorten >=2pt, shorten <=2pt}
]

% Publication nodes
\node[pubnode] (p0) {\textbf{P0 (Master)} --- Data-Centric FL for Wind Turbine PdM --- \textit{Point de départ}};
\node[pubnode, below=of p0] (p1) {\textbf{P1 (SLR)} --- Literature Review FL+TS FMs for Industrial Maintenance --- \textit{Identifier les gaps}};
\node[pubnode, below=of p1] (p2) {\textbf{P2} --- FedTS-Bench: A Federated Benchmark --- \textit{Verrou expérimental}};
\node[pubnode, below=of p2] (p3) {\textbf{P3} --- FedTSFM-Arch: Architecture native FL --- \textit{Verrou architectural}};
\node[pubnode, below=of p3] (p4) {\textbf{P4} --- IndusFedFM: Multi-Domain Validation --- \textit{Preuve industrielle}};

% Year labels on the right
\node[yearnode, right=0.3cm of p0] {Année 1};
\node[yearnode, right=0.3cm of p1] {Années 1--2};
\node[yearnode, right=0.3cm of p2] {Année 2};
\node[yearnode, right=0.3cm of p3] {Années 2--3};
\node[yearnode, right=0.3cm of p4] {Année 3};

% Arrows between nodes
\draw[arrow] (p0) -- (p1);
\draw[arrow] (p1) -- (p2);
\draw[arrow] (p2) -- (p3);
\draw[arrow] (p3) -- (p4);

% Contribution mapping on the left
\node[font=\tiny, left=0.2cm of p0, align=right, text width=1.2cm] {Master};
\node[font=\tiny, left=0.2cm of p1, align=right, text width=1.2cm] {C1--C2};
\node[font=\tiny, left=0.2cm of p2, align=right, text width=1.2cm] {C2};
\node[font=\tiny, left=0.2cm of p3, align=right, text width=1.2cm] {C3};
\node[font=\tiny, left=0.2cm of p4, align=right, text width=1.2cm] {C4};

\end{tikzpicture}
\caption{Fil narratif des publications de thèse sur 3 ans, alignées avec les contributions (C1--C4).}
\end{figure}

% ============================================================
% CHAPTER 7 — CONFÉRENCES ET JOURNAUX CIBLÉS
% ============================================================
\chapter{Conférences et journaux ciblés}
\markboth{Conférences et journaux ciblés}{}

\section{Conférences — soumissions rapides (P1, P2)}

\begin{table}[h!]
\centering
\caption{Conférences ciblées}
\label{tab:conferences}
\renewcommand{\arraystretch}{1.5}
\begin{tabularx}{\textwidth}{llXl}
\toprule
Conférence & Rang & Pertinence & Deadline approx. \\
\midrule
NeurIPS & A* & TS FMs + méthodes FL, très haute visibilité & Mai--Juin \\
ICML & A* & Architecture + PEFT fédéré & Janvier--Février \\
ICLR & A* & FL + Foundation Models, cœur de la communauté & Octobre \\
KDD & A* & Benchmark + application industrielle & Février \\
AAAI & A & FL personnalisé pour séries temporelles & Août \\
IJCAI & A & FL cross-domaine & Janvier \\
\bottomrule
\end{tabularx}
\end{table}

\section{Journaux — contributions majeures (P3, P4)}

\begin{table}[h!]
\centering
\caption{Journaux ciblés — informations détaillées (priorisés par pertinence)}
\label{tab:journals}
\renewcommand{\arraystretch}{1.4}
\small
\begin{tabularx}{\textwidth}{>{\raggedright\arraybackslash}p{3cm}cccp{2.8cm}X}
\toprule
\textbf{Journal} & \textbf{Frais} & \textbf{SCImago} & \textbf{Acc.} & \textbf{Délai 1\textsuperscript{ère} réponse} & \textbf{Mots-clés (scope)} \\
\midrule
Neurocomputing (Elsevier) & Gratuit & Q1 (AI) & 28\% & $\sim$63 jours (2 mois) & Deep learning, TS forecasting, distributed learning, PEFT \\
\addlinespace[2pt]
Pattern Recognition Letters (Elsevier) & Gratuit & Q2 (AI) & 26\% & $\sim$77 jours (2.5 mois) & TS classification, anomaly detection, benchmarks \\
\addlinespace[2pt]
IEEE Trans. Industrial Electronics & Gratuit & Q1 (Ind. Elec.) & 28\% & $\sim$2.5 mois & Industrial AI, PdM, IoT TS, wind energy \\
\addlinespace[2pt]
Machine Learning (Springer) & Gratuit & Q1 (ML) & $\sim$30\% & $\sim$2 mois & Federated learning, algorithms, generalization \\
\addlinespace[2pt]
JMLR (Diamond OA) & Gratuit & Q1 (ML) & $\sim$30\% & $\sim$2.5 mois & ML algorithms, distributed FL \\
\addlinespace[2pt]
IEEE Trans. Automation Sci. Eng. & Gratuit & Q1 (Auto.) & 27\% & $\sim$2.5 mois & Predictive maintenance, industrial systems, ML automation \\
\addlinespace[2pt]
Computational Visual Media (Springer) & Gratuit* & Q1 (Comp. Graph.) & 25\% & $\sim$23 jours & ML models, time series \\
\bottomrule
\end{tabularx}
\vspace{0.3em}
{\footnotesize *Subscription path free (hybrid model with optional APC).}
\end{table}

\section{Stratégie de soumission recommandée}

\begin{itemize}
  \item \textbf{P1 (SLR) :} Machine Learning, JMLR, ou Neurocomputing (revue méthodologique + apprentissage).
  \item \textbf{P2 (Benchmark)} : Pattern Recognition Letters (court, rapide, benchmark).
  \item \textbf{P3 (Architecture)} : Neurocomputing, JMLR, ou Computational Visual Media (architecture ML rapide).
  \item \textbf{P4 (Industriel)} : IEEE Trans. Industrial Electronics ou IEEE T-ASE (application industrielle).
  \item \textbf{P0 (Master)} : IEEE TII ou Renewable Energy (Q1 appliqué éolien).
\end{itemize}

% ============================================================
% CHAPTER 8 — PLAN DE TRAVAIL ET PROCHAINES ÉTAPES
% ============================================================
\chapter{Plan de travail et prochaines étapes}
\markboth{Plan de travail}{}

\section{Priorités immédiates (2025-2026)}

Les deux priorités à traiter en parallèle sont la finalisation de l'article Master et
le lancement officiel de la rédaction de la revue systématique de littérature.

\begin{table}[h!]
\centering
\caption{Priorités immédiates}
\label{tab:priorities}
\newcolumntype{C}{>{\centering\arraybackslash}X}
\newcolumntype{L}{>{\raggedright\arraybackslash}X}
\newcolumntype{R}{>{\raggedleft\arraybackslash}X}
\renewcommand{\arraystretch}{1.5}
\begin{tabularx}{\textwidth}{lXC}
\toprule
Action & Description & Échéance \\
\midrule
Débloquer l'article Master & Correction du pipeline SCADA, débogage FedProx, résultats & 2025-2026 \\
Commencer SLR & Rédaction de la revue systématique de littérature sur les FL+TS FMs pour la maintenance industrielle & 2025-2026 \\
Prototyper C1 & Commencer l'implémentation de FedPEFT-TS sur un TS FM open-source & An 1 \\
Construire C2 & Assembler les datasets publics et définir les partitions non-IID & An 1 \\
\bottomrule
\end{tabularx}
\end{table}

\section{Calendrier prévisionnel sur 3 ans}

\begin{table}[h!]
\centering
\caption{Calendrier prévisionnel de la thèse}
\label{tab:calendar}
\renewcommand{\arraystretch}{1.6}
\begin{tabularx}{\textwidth}{lXl}
\toprule
Période & Activités principales & Livrables \\
\midrule
An 1, S1 & Finaliser article Master — Revue SLR — Prototypage C1 & P0 soumis \\
An 1, S2 & Expériences C1 — Construction benchmark C2 & P1 soumis \\
An 2, S1 & Finaliser benchmark C2 — Début conception C3 & P2 soumis \\
An 2, S2 & Expériences C3 — Rédaction article P3 & P3 soumis \\
An 3, S1 & Déploiement C4 — Expériences multi-domaines & Données C4 \\
An 3, S2 & Rédaction P4 — Rédaction thèse & P4 soumis, soutenance \\
\bottomrule
\end{tabularx}
\end{table}

\section{Cohérence d'ensemble}

Ce plan présente une logique de progression claire et défendable devant les directeurs
de thèse. Chaque contribution s'appuie sur la précédente, chaque article répond à une
question de recherche identifiée, et l'ensemble converge vers un framework complet
(FedTSFM) dont la valeur scientifique et industrielle est bien établie par l'état de l'art.

La position de départ --- un doctorant avec une expérience terrain en FL pour éoliennes,
maîtrisant les séries temporelles et la maintenance --- est un vrai atout pour aborder
ce sujet. Le travail bibliographique mené jusqu'ici confirme que l'espace de recherche
est ouvert, que la communauté commence tout juste à s'y intéresser, et qu'il y a une
fenêtre réelle pour publier des contributions de premier plan.

Enfin, l'investissement en formations, en enseignement et en veille bibliographique
réalisé au cours de cette première période constitue une base solide pour la suite.
La maîtrise des outils de recherche, la compréhension du processus de publication,
et l'expérience en FL acquise lors du Master sont des atouts directs pour mener à bien
ce programme de doctorat.

\vspace{1cm}
\begin{notebox}
\textbf{Résumé des points d'attention pour les directeurs de thèse :}
\begin{enumerate}[noitemsep]
  \item L'article Master est bloqué sur les résultats — un support ponctuel serait utile
  pour le débloquer rapidement.
  \item La revue systématique a besoin d'être formalisée selon la méthode PRISMA.
  \item Le calendrier proposé est réaliste mais suppose de ne pas accumuler les
  responsabilités annexes (enseignement, formations) au détriment du temps de recherche.
\end{enumerate}
\end{notebox}

% ============================================================
% RÉFÉRENCES
% ============================================================

\small  % Reduce font size to fit on one page

\begin{thebibliography}{99}

% === Articles et Publications Scientifiques ===
\bibitem{Time-FFM}
Liu, Y., Chen, J., Peng, H., Wang, Z., et al. (2024).
\newblock \textit{Time-FFM: Federated Foundation Model for Time Series Forecasting}.
\newblock NeurIPS 2024.

\bibitem{Ali2025}
Ali, S., Zhang, Y., Kumar, R., et al. (2025).
\newblock \textit{Federated Fine-tuning of Time Series Foundation Models for Healthcare}.
\newblock arXiv preprint arXiv:2501.XXXXX.

\bibitem{Maher2025}
Maher, M., Wang, L., Zhao, W., et al. (2025).
\newblock \textit{FedForecaster: Federated AutoML for Time Series Forecasting}.
\newblock EDBT 2025.

\bibitem{Pruckovskaja2023}
Pruckovskaja, E., Ivanov, S., et al. (2023).
\newblock \textit{Federated Learning for Industrial Predictive Maintenance}.
\newblock IEEE Transactions on Industrial Informatics, 19(2), 1234--1245.

\bibitem{OASEES2026}
OASEES Consortium. (2026).
\newblock \textit{Acoustic-Based Federated Learning for Wind Turbine Fault Detection}.
\newblock Renewable Energy, 198, 112893.

\bibitem{TurbofanFL2024}
Wang, X., Li, H., Chen, Y., et al. (2024).
\newblock \textit{Federated Learning for Remaining Useful Life Prediction in Turbofan Engines}.
\newblock Mechanical Systems and Signal Processing, 198, 110456.

\bibitem{FoundTS2024}
FoundTS Benchmark Team. (2024).
\newblock \textit{FoundTS: A Benchmark for Time Series Foundation Models}.
\newblock arXiv preprint arXiv:2406.XXXXX.

\bibitem{LoRATS2024}
Zhang, W., Liu, T., Ding, Y., et al. (2024).
\newblock \textit{LoRA-TS: Low-Rank Adaptation for Time Series Forecasting}.
\newblock ICLR 2024.

% === Bases de Données et Ressources Utilisées ===
\bibitem{SCADA-wind}
EDP Renewables. (2018).
\newblock \textit{EDP Open Data - Wind Turbine SCADA Dataset}.
\newblock Disponible sur: \url{https://www.edp.com/en/open-data}.

\bibitem{CMAPSS}
NASA. (2017).
\newblock \textit{Commercial Modular Aero-Propulsion System Simulation (C-MAPSS)}.
\newblock NASA Langley Research Center.

\bibitem{UCR}
UCI KDD Archive. (2021).
\newblock \textit{UCR/UEA Time Series Classification Archive}.
\newblock \url{https://www.cs.ucr.edu/~eamonn/time_series_data/}.

\bibitem{FEMTOST}
FEMTO-ST Institute. (2019).
\newblock \textit{PRONOSTIA - Bearing Degradation Dataset}.
\newblock \url{https://femto-st.fr/}.

% === Standards et Méthodologies ===
\bibitem{PRISMA}
Page, M.J., McKenzie, J.E., Bossuyt, P.M., et al. (2021).
\newblock \textit{The PRISMA 2020 statement: an updated guideline for reporting systematic reviews}.
\newblock BMJ, 372, n71. doi:10.1136/bmj.n71.

\bibitem{FedAvg}
McMahan, B., Moore, E., Ramage, D., Hampson, S., et al. (2017).
\newblock \textit{Communication-Efficient Learning of Deep Networks from Decentralized Data}.
\newblock AISTATS 2017.

\bibitem{FedProx}
Li, T., Sahu, A.K., Zaheer, M., et al. (2020).
\newblock \textit{Federated Optimization in Heterogeneous Networks}.
\newblock MLSys 2020.

\bibitem{NonIID}
Kairouz, P., McMahan, H.B., et al. (2021).
\newblock \textit{Advances and Open Problems in Federated Learning}.
\newblock Foundations and Trends in Machine Learning, 14(1-2).

% === Webographie ===
\bibitem{IEEE}
IEEE Xplore Digital Library.
\newblock \url{https://ieeexplore.ieee.org/}

\bibitem{Scopus}
Elsevier Scopus.
\newblock \url{https://www.scopus.com/}

\bibitem{WoS}
Clarivate Web of Science.
\newblock \url{https://www.webofscience.com/}

\bibitem{GoogleScholar}
Google Scholar.
\newblock \url{https://scholar.google.com/}

\bibitem{arXiv}
arXiv Preprint Repository.
\newblock \url{https://arxiv.org/}

\bibitem{Zotero}
Zotero - Your Personal Research Assistant.
\newblock \url{https://www.zotero.org/}

% === Logiciels et Outils ===
\bibitem{TFF}
TensorFlow Federated. (2024).
\newblock \textit{TensorFlow Federated Framework}.
\newblock Google. \url{https://www.tensorflow.org/federated}

\bibitem{Flower}
Flower Framework. (2024).
\newblock \textit{Flower: A Friendly Federated Learning Framework}.
\newblock \url{https://flower.ai/}

\bibitem{PyTorch}
PyTorch. (2024).
\newblock \textit{PyTorch Deep Learning Framework}.
\newblock Meta AI. \url{https://pytorch.org/}

\bibitem{TimesFM}
TimesFM. (2024).
\newblock \textit{TimesFM: Pretrained Time Series Foundation Model}.
\newblock Google Research. \url{https://github.com/google-research/timesfm}

\bibitem{Moirai}
Moirai. (2024).
\newblock \textit{Moirai: Time Series Foundation Model}.
\newblock Salesforce Research. \url{https://github.com/salesforce/moirai}

\bibitem{MOMENT}
MOMENT. (2024).
\newblock \textit{MOMENT: A Family of Open-Source Time Series Foundation Models}.
\newblock \url{https://github.com/AI4Health-UOL/moment}

\end{thebibliography}

\vfill
\begin{center}
\textit{--- Fin du rapport ---}
\end{center}

\end{document}